\begin{table}[htbp]
\centering
\caption{Summary Statistics}
\label{tab:summary}
\begin{threeparttable}

\textbf{Panel A: Sample Counts}
\medskip

\begin{tabular}{lccc}
\toprule
 & Full Sample & \makecell{Drafted Years} & 1944 \\
\midrule
Unique workers & 92,439 & 54,254 & 9,323 \\
Unique kyoku & 119 & 48 & 8 \\
Unique ka (ka $\times$ kyoku) & 4,967 & 1,464 & 276 \\
Unique kakari (kakari $\times$ ka $\times$ kyoku) & 9,234 & 1,690 & 383 \\
Unique women & 5,968 & 3,720 & 702 \\
Workers at rank 3 (shiji/gishi) & 15,296 & 4,989 & 470 \\
Workers at rank 2 (shoki/gite) & 124,127 & 74,352 & 8,013 \\
Workers at rank 1 (yatoi/shokutaku) & 99,555 & 29,312 & 789 \\
\bottomrule
\end{tabular}

\bigskip

\textbf{Panel B: Kakari-Level Averages}
\medskip

\begin{tabular}{lccc}
\toprule
 & Full Sample & \makecell{Drafted Years} & 1944 \\
\midrule
Workers per kakari & 22.7 (102.4) & 58.3 (206.6) & 24.3 (38.7) \\
Unique positions per kakari & 2.2 (1.7) & 3.4 (2.6) & 3.4 (2.3) \\
Workers per rank-3 role & 3.5 (22.5) & 5.3 (47.2) & 2.3 (3.8) \\
Workers per rank-2 role & 12.8 (48.5) & 19.3 (66.1) & 8.5 (14.7) \\
Workers per rank-1 role & 11.5 (36.1) & 23.6 (67.4) & 6.3 (8.3) \\
\bottomrule
\end{tabular}

\begin{tablenotes}[flushleft]
\footnotesize
\item \textit{Notes:} Panel~A reports unique counts for the full sample, the subset of years with any drafted workers (1938, 1940, 1941, 1943, 1944), and 1944 alone. Rank 3 = \emph{shiji}/\emph{gishi} (supervisory); Rank 2 = \emph{shoki}/\emph{gite} (clerical/technical staff); Rank 1 = \emph{yatoi}/\emph{shokutaku} (temporary/contract). Panel~B reports kakari-level averages (standard deviations in parentheses). A kakari is the lowest organizational unit (team), nested within ka (section) and kyoku (bureau).
\end{tablenotes}
\end{threeparttable}
\end{table}
